\section{SmallWood Programming Model}
\label{sec:smallwood-model}

This section records the semantic SmallWood statements used by ARC-SPRUCE.  The concrete row layouts,
DECS/LVCS/PCS openings, and proof-stack details are deferred to Appendix~\ref{app:smallwood}.  The prover
commits to packed witness rows over a set $\Omega$ and proves PACS-style constraints on those committed rows.

\subsection{Proof interface}
Fix a packing set $\Omega=\{\omega_1,\ldots,\omega_s\}\subset\Fq$ and a disjoint hiding set
$\Omega'\subset\Fq\setminus\Omega$.  A compiled witness is a matrix of row values over $\Fq$.  Each row is
interpolated to a low-degree polynomial whose values on $\Omega$ encode the witness and whose values on
$\Omega'$ are random masks.  SmallWood then checks parallel constraints pointwise over $\Omega$ and, when
needed, aggregated constraints through $\Omega$-sum identities.

\subsection{Issuance as a SmallWood statement}
\label{subsec:presign-compiled}
The public issuance statement is
\[
  (\ComPublicParam,\vecc,T,r_0^I,r_1^I,\mA,B),
\]
and the hidden witness is
\[
  (\vecm,\veck,r_0^H,r_1^H,\bar r,r_0,r_1,Z).
\]
Semantically, the proof establishes four claims.
First, the commitment opens to the same message and holder-side randomness used later:
\[
  \vecc=\mACom[\vecm\Vert\veck\Vert r_0^H\Vert r_1^H\Vert\bar r]^\top.
\]
Second, the public issuer randomness is combined correctly with the committed holder randomness:
\[
  r_0=\centerfunc(r_0^H+r_0^I),\qquad
  r_1=\centerfunc(r_1^H+r_1^I).
\]
Third, the inverse witness is valid:
\[
  (B_3-r_1)\Had Z=1.
\]
Fourth, the public target is the BB-tran target for the same hidden message and centered randomness:
\[
  T=B_0+B_1(\vecm\Vert\veck)+B_2r_0+Z.
\]
All message and randomness rows are constrained to lie in the coefficient interval
$[-\BoundMsg,\BoundMsg]$.

The corresponding committed rows consist of carrier rows for the packed message, the holder-side
randomness, commitment-only randomness, centered randomness, carry rows for the centering map, and the
inverse witness $Z$.  The pre-sign proof is dominated by commitment linearity, centering, range membership,
the inverse-witness equation, and the target identity above.

\subsection{Showing as a SmallWood statement}
\label{subsec:show-compiled}
The public showing statement is
\[
  (\nonce,\prftag,\mA,B,\BoundMsg,\BoundSig)
\]
with verifier state handled outside the proof.  The hidden witness is
\[
  (\vecu,\vecm,\veck,r_0,r_1,Z).
\]
The proof certifies
\begin{align*}
  &(B_3-r_1)\Had Z=1,\\
  &\mA\vecu=B_0+B_1(\vecm\Vert\veck)+B_2r_0+Z,\\
  &\prftag=\PRF(\veck,\nonce),
\end{align*}
with coefficient bounds on $\vecm,\veck,r_0,r_1$ and the public shortness bound on $\vecu$.

In the compiled statement, the signature surface proves shortness of $\vecu$ and binds the derived target to
$\mA\vecu$.  The BB-tran part carries source rows for $T=\mA\vecu$, $\vecm\Vert\veck$, $r_0$, $r_1$, and
$Z$.  The PRF companion rows bind the same hidden key $\veck$ to the public tag.

\subsection{PRF companion relation}
\label{subsec:prf-checkpointed}
The PRF relation ties the public tag to the same hidden key that appears in the signed message.  Selector
constraints bind the PRF key lanes to $\veck$; checkpoint constraints certify the nonlinear S-box steps of the
Poseidon2-like permutation; and the final feed-forward and truncation constraints bind the resulting lanes to
$(\nonce,\prftag)$.

\subsection{Statement strength}
\label{subsec:statement-strength}
The showing statement certified by SmallWood is the coefficient-domain BB-tran identity in witness form:
\[
  (B_3-r_1)\Had Z=1,
  \qquad
  \mA\vecu=B_0+B_1(\vecm\Vert\veck)+B_2r_0+Z.
\]
Equivalently, it implies the eliminated identity
\[
  (B_3-r_1)\Had\bigl(\mA\vecu-B_0-B_1(\vecm\Vert\veck)-B_2r_0\bigr)=1.
\]
Thus any accepting showing proof yields, by knowledge soundness, a witness satisfying the abstract BB-tran
signature relation and the PRF-tag relation.

\begin{lemma}[Witness consequence of an accepting showing proof]\label{lem:showing-witness-consequence}
If the post-sign NIZK is knowledge sound for the compiled showing relation, then every accepting proof yields
a witness $(\vecu,\vecm,\veck,r_0,r_1,Z)$ satisfying
\[
  \mA\vecu=h_{\mathrm{tran},\vecm\Vert\veck,(r_0,r_1)}(B)
\]
and $\prftag=\PRF(\veck,\nonce)$.
\end{lemma}

\begin{proof}
The extractor returns a witness satisfying the compiled equations above.  The inverse equation gives
$Z=(B_3-r_1)^{-1}$ under admissibility, and substituting it into the target equation gives the abstract
BB-tran relation.
\end{proof}
